\textbf{1.} On a $1! = 1 < 2$, $2! = 2 < 4$, $3! = 6 < 8$, mais
$4! = 24 > 16 = 2^4$. L'entier $n = 4$ est un contre-exemple : la propriété
n'est pas vraie pour tout $n \geqslant 1$.

\textbf{2.} L'inégalité $n! < 3^n$ est vraie pour $n = 1, \ldots, 6$
(par exemple $6! = 720 < 729 = 3^6$), mais $7! = 5040 > 2187 = 3^7$.
Elle est donc fausse en général.

\textbf{3.} L'inégalité $n! < 5^n$ résiste plus longtemps :
$11! = 39\,916\,800 < 48\,828\,125 = 5^{11}$, mais
$12! = 479\,001\,600 > 244\,140\,625 = 5^{12}$. Elle est donc fausse : il
suffit d'\emph{un seul} contre-exemple, $n = 12$.

\emph{Remarque.} Une fois l'inégalité renversée, elle le reste : si
$n! \geqslant a^n$ et $n + 1 > a$, alors $(n+1)! = (n+1)\,n! > a \cdot a^n = a^{n+1}$.
Ainsi $n! > 2^n$ pour tout $n \geqslant 4$, $n! > 3^n$ pour tout $n \geqslant 7$
et $n! > 5^n$ pour tout $n \geqslant 12$. Plus les premiers cas favorables
sont nombreux, plus il est tentant de généraliser à tort. On notera enfin que,
si l'on autorisait $n = 0$, il serait aussi un contre-exemple aux trois
inégalités, puisque $0! = 1 = a^0$ : les cas limites sont de bons candidats.
\hfill $\blacksquare$
